% ============================================================================
% WORKBOOK — Section 4.6: Rewriting Expressions to Solve Problems
% CCSS 7.EE.A.2
% Practice-focused reinforcement for the study guide topic
% ============================================================================
\section{Rewriting Expressions to Solve Problems}
\topicTitle{Rewriting Expressions to Solve Problems}

\begin{quickReview}{Rewriting Expressions}
The \textbf{same} situation can be written in \textbf{different} forms. Each form reveals something useful.

\begin{itemize}[leftmargin=*]
    \item $a + 0.05a = 1.05a$ $\to$ ``add $5\%$'' is the same as ``multiply by $1.05$''
    \item $p - 0.20p = 0.80p$ $\to$ ``take off $20\%$'' is the same as ``pay $80\%$''
\end{itemize}

\medskip
\textbf{Why rewrite?} One form may be easier to compute, or it may highlight a relationship you did not see at first.

\smallskip
\textbf{Example:} A $\$60$ jacket is $25\%$ off. \quad Form 1: $60 - 0.25(60) = 60 - 15 = 45$. \quad Form 2: $0.75(60) = 45$. Both work — Form 2 is faster.
\end{quickReview}

% ── Warm-Up ──────────────────────────────────────────────────────────────────
\begin{practiceBox}[funGreen]{\faSun~Warm-Up}

\practiceHeader[funGreen]{Write each as a single multiplication.}

\begin{multicols}{2}
\prob $x + 0.10x$ \answerBlank[3cm]
\answerExplain{$1.10x$}{Factor out $x$: $x(1 + 0.10) = 1.10x$. This represents a $10\%$ increase.}

\prob $p - 0.25p$ \answerBlank[3cm]
\answerExplain{$0.75p$}{Factor out $p$: $p(1 - 0.25) = 0.75p$. This represents a $25\%$ decrease.}

\prob $n + 0.50n$ \answerBlank[3cm]
\answerExplain{$1.50n$}{Factor out $n$: $n(1 + 0.50) = 1.50n$. This represents a $50\%$ increase.}

\prob $c - 0.15c$ \answerBlank[3cm]
\answerExplain{$0.85c$}{Factor out $c$: $c(1 - 0.15) = 0.85c$. This represents a $15\%$ decrease.}

\prob $y + 0.06y$ \answerBlank[3cm]
\answerExplain{$1.06y$}{Factor out $y$: $y(1 + 0.06) = 1.06y$. This represents a $6\%$ increase, like adding sales tax.}

\prob $d - 0.40d$ \answerBlank[3cm]
\answerExplain{$0.60d$}{Factor out $d$: $d(1 - 0.40) = 0.60d$. This represents a $40\%$ discount — you pay $60\%$.}
\end{multicols}
\end{practiceBox}

% ── Practice A: Interpreting Equivalent Forms ────────────────────────────────
\begin{practiceBox}{What Does Each Form Mean?}

\practiceHeader{Rewrite and interpret.}

\prob A shirt costs $\$s$. Sales tax is $8\%$. Write two equivalent expressions for the total cost.
\answerBlank[5cm]
\answerExplain{$s + 0.08s$ and $1.08s$}{The tax amount is $0.08s$. Add it to the price: $s + 0.08s$. Factor: $s(1 + 0.08) = 1.08s$. Both represent the total cost including tax.}

\prob A store marks up items by $30\%$. If an item costs the store $\$c$, write the selling price as a single term.
\answerBlank[3cm]
\answerExplain{$1.30c$}{Markup is $0.30c$. Selling price: $c + 0.30c = 1.30c$. The store sells the item for $130\%$ of its cost.}

\prob A coupon gives $\$5$ off, then an additional $10\%$ off the reduced price. Write an expression for the final price of an item that originally costs $\$p$.
\answerBlank[5cm]
\answerExplain{$0.90(p - 5)$ or $0.9p - 4.5$}{First subtract $\$5$: $p - 5$. Then take $10\%$ off the result: $0.90(p - 5)$. Expanding: $0.9p - 4.5$.}

\prob Rewrite $2(3x + 4) - x$ in simplest form. \answerBlank[3cm]
\answerExplain{$5x + 8$}{Expand: $6x + 8 - x$. Combine like terms: $(6x - x) + 8 = 5x + 8$.}

\end{practiceBox}

% ── Practice B: Simplify and Interpret ───────────────────────────────────────
\begin{practiceBox}[funTeal]{Simplify to Reveal the Pattern}

\prob Show that ``increase by $20\%$'' is the same as ``multiply by $1.20$'' using the variable $a$.
\answerBlank[5cm]
\answerExplain{$a + 0.20a = a(1 + 0.20) = 1.20a$}{Start with $a$, add $20\%$: $a + 0.20a$. Factor out $a$: $a(1 + 0.20) = 1.20a$. Both are equivalent, confirming the statement.}

\prob Simplify $4(x + 3) - 2(x - 1)$. What do you notice? \answerBlank[4cm]
\answerExplain{$2x + 14$}{Expand: $4x + 12 - 2x + 2$. Combine: $(4x - 2x) + (12 + 2) = 2x + 14$. The expression simplifies to a linear expression with a positive constant.}

\prob A babysitter charges $\$12$ per hour plus a $\$5$ travel fee. She babysat for $h$ hours on Friday and $h + 2$ hours on Saturday. Write a simplified expression for her total earnings.
\answerBlank[5cm]
\answerExplain{$24h + 34$}{Friday: $12h + 5$. Saturday: $12(h + 2) + 5 = 12h + 24 + 5 = 12h + 29$. Total: $(12h + 5) + (12h + 29) = 24h + 34$.}

\prob A cell phone plan costs $\$40$ per month plus $\$0.05$ per text. Another plan costs $\$25$ plus $\$0.15$ per text. Write a simplified expression for how much more Plan 1 costs than Plan 2.
\answerBlank[5cm]
\answerExplain{$15 - 0.10t$}{Plan 1: $40 + 0.05t$. Plan 2: $25 + 0.15t$. Difference: $(40 + 0.05t) - (25 + 0.15t) = 40 + 0.05t - 25 - 0.15t = 15 - 0.10t$. Plan 1 costs more when $t < 150$ texts.}

\end{practiceBox}

% ── Practice C: True or False ────────────────────────────────────────────────
\begin{practiceBox}[funPurple]{True or False?}

\trueOrFalse{$x + 0.07x$ and $1.07x$ are equivalent expressions.}
\answerExplain{True}{Factor out $x$: $x(1 + 0.07) = 1.07x$. The two expressions always have the same value for every value of $x$.}

\trueOrFalse{A $30\%$ discount followed by a $30\%$ increase brings you back to the original price.}
\answerExplain{False}{Start with price $p$. After a $30\%$ discount: $0.70p$. After a $30\%$ increase: $1.30(0.70p) = 0.91p$. The final price is only $91\%$ of the original — not $100\%$.}

\trueOrFalse{$5(x - 2)$ and $5x - 10$ represent the same quantity.}
\answerExplain{True}{Expand using the distributive property: $5(x - 2) = 5x - 10$. They are equivalent expressions.}

\end{practiceBox}

% ── Practice D: Word Problems ────────────────────────────────────────────────
\begin{practiceBox}[funBlue]{\faGlobe~Real-World Rewriting}

\wordProblem{A jacket costs $\$80$ and is on sale for $35\%$ off. Use the expression $0.65(80)$ to find the sale price.}{dollars}
\answerExplain{$\$52$}{A $35\%$ discount means you pay $100\% - 35\% = 65\%$ of the price. Multiply: $0.65 \times 80 = 52$.}

\wordProblem{A car's value depreciates $12\%$ each year. If it is worth $\$v$ now, write an expression for its value after one year.}{dollars}
\answerExplain{$0.88v$}{Depreciates $12\%$ means the car keeps $100\% - 12\% = 88\%$ of its value. Expression: $v - 0.12v = 0.88v$.}

\wordProblem{Two friends split the cost of a pizza that costs $\$p$ plus $6\%$ tax equally. Write and simplify an expression for each person's share.}{dollars}
\answerExplain{$0.53p$}{Total with tax: $p + 0.06p = 1.06p$. Split equally: $\frac{1.06p}{2} = 0.53p$. Each person pays $\$0.53p$.}

\end{practiceBox}

% ── Challenge ────────────────────────────────────────────────────────────────
\begin{challengeBox}
\prob A store increases prices by $10\%$, then offers a $10\%$ sale. A customer says, ``The sale just cancels the increase, so I pay the original price.'' Show mathematically whether the customer is correct.
\answerBlank[6cm]
\answerExplain{The customer is wrong; final price is $0.99p$}{Start with price $p$. Increase by $10\%$: $1.10p$. Decrease by $10\%$: $0.90(1.10p) = 0.99p$. The final price is $99\%$ of the original — $1\%$ less, not the same.}

\prob Show that ``buy $3$, get $1$ free'' on items costing $\$d$ each is a $25\%$ discount per item.
\answerBlank[6cm]
\answerExplain{True — price per item is $0.75d$}{You pay for $3$ items but receive $4$: total cost $= 3d$ for $4$ items. Cost per item $= \frac{3d}{4} = 0.75d$. Since $0.75d = d - 0.25d$, this is a $25\%$ discount on each item.}
\end{challengeBox}

\encouragement{Outstanding work! Rewriting expressions is a superpower that makes hard problems much easier!}
